% ---------------------------------------------------------------------------
\section{Applications}
\label{sec:applications}
% ---------------------------------------------------------------------------

In this section, we exploit the directional consistency of ACE's gradient with respect to HEVC's bitrate to perform compression-aware optimization to defend against adversarial attacks at inference time (\S~\ref{sec:adv_pur}), improve machine learning model robustness through training-time interventions (\S~\ref{sec:aug}), improve attack success rate against compression-based defenses (\S~\ref{sec:survival}), facilitate effective split-inference (\S~\ref{sec:split}), and guide diffusion to high-quality yet compressible outputs (\S~\ref{sec:diffguide}).
In each case, ACE is used as is, with no requirement for training or hyperparameter tuning.

\subsection{Adversarial Purification}\label{sec:adv_pur}

Adversarial perturbations are crafted to maximize a classifier's error while remaining small in the $\ell_\infty$ norm.
Here, we demonstrate that ACE's gradients can be used as a purification signal to reverse the adversarial nature of the perturbations.

Given an adversarial image $\bx_\mathrm{adv}$, we run Adam~\citep{kingma2014adam} for $T$ steps minimizing ACE with a fidelity constraint
\begin{equation}
  \bx_{t+1} = \mathrm{clip}_{[0,1]}\!\left(
    \bx_t - \eta\nabla_{\bx_t}
    \left(\frac{1}{2}\|\bx_t - \bx_\mathrm{adv}\|_2^2
          + \lambda\,\hat{r}(\bx_t)\right)
  \right).
  \label{eq:purify}
\end{equation}
The $\ell_2$ term keeps the purified image close to the adversarial input, preventing the proxy from smoothing away semantically relevant content.
The hyperparameters $T$ and $\lambda$ trade off compute budget and fidelity, respectively.

Figure~\ref{fig:purification} summarizes performance on ResNet-50~\citep{he2016resnet} trained on ImageNet.
The left panel reports top-1 accuracy at $T=10$ across four attack configurations, and shows that through ACE's gradient \eqref{eq:purify} can help recover $32\%$--$45\%$ of the clean--adversarial gap, compared to $3\%$--$20\%$ for JPEG compression at quality $Q$ equal to $75$---the non-differentiable hard-quantization analog of the proxy.
The right panel traces the operating envelope for all four attacks as $T$ increases from $0$ to $30$, showing that recovery peaks at $T=10$ for every configuration, and that subsequent steps erode clean accuracy without additional adversarial benefit.

\begin{figure}[ht]
  \centering
  \begin{minipage}[t]{0.46\linewidth}
    \centering
    \vspace{0pt}
    \scriptsize
    \begin{tabular}{lcccc}
      \toprule
      Attack & $\varepsilon$ & Adv & JPEG & DCT \\
      \midrule
      FGSM   & 4/255 & 6.3 & 17.0 (15.6\%) & 33.2 (38.8\%) \\
      FGSM   & 8/255 & 6.9 & 13.4 (9.4\%) & 29.2 (32.3\%) \\
      PGD-10 & 4/255 & 0.0 & 15.2 (20.1\%) & 34.2 (45.2\%) \\
      PGD-10 & 8/255 & 0.0 & 2.1 (2.7\%) & 30.0 (39.6\%) \\
      \bottomrule
    \end{tabular}
  \end{minipage}%
  \hfill
  \begin{minipage}[t]{0.51\linewidth}
    \centering
    \vspace{0pt}
    \begin{tikzpicture}
    \begin{axis}[
      width  = \linewidth,
      height = 0.65\linewidth,
      xlabel = {Clean accuracy drop (\%)},
      ylabel = {Gap recovery (\%)},
      xmin=-2,  xmax=60,
      ymin=-3,  ymax=52,
      xtick={0,10,20,30,40,50},
      ytick={0,10,20,30,40,50},
      tick label style={font=\scriptsize},
      label style={font=\small},
      legend style={font=\tiny, at={(0.98,0.02)}, anchor=south east,
                    row sep=-2pt},
      legend columns=2,
    ]
      \addplot[color=proxyblue, line width=1.8pt,
               mark=*, mark size=1.8pt, mark options={fill=proxyblue}]
        table[col sep=comma, x=clean_drop_pp, y=recovery_pct]
        {figures/data/lam010_pgd_eps4.csv};
      \addlegendentry{PGD-10 $\varepsilon{=}4$}

      \addplot[color=proxyblue, line width=1.8pt, dashed,
               mark=square*, mark size=1.6pt, mark options={fill=proxyblue}]
        table[col sep=comma, x=clean_drop_pp, y=recovery_pct]
        {figures/data/lam010_pgd_eps8.csv};
      \addlegendentry{PGD-10 $\varepsilon{=}8$}

      \addplot[color=black!50, line width=1.8pt,
               mark=triangle*, mark size=1.8pt, mark options={fill=black!50}]
        table[col sep=comma, x=clean_drop_pp, y=recovery_pct]
        {figures/data/lam010_fgsm_eps4.csv};
      \addlegendentry{FGSM $\varepsilon{=}4$}

      \addplot[color=black!50, line width=1.8pt, dashed,
               mark=diamond*, mark size=1.8pt, mark options={fill=black!50}]
        table[col sep=comma, x=clean_drop_pp, y=recovery_pct]
        {figures/data/lam010_fgsm_eps8.csv};
      \addlegendentry{FGSM $\varepsilon{=}8$}

    \end{axis}
    \end{tikzpicture}
  \end{minipage}
  \caption{
  \textbf{ACE's gradient signals can purify adversarial perturbations to recover classifier performance.}
  We consider purification using \eqref{eq:purify} with $\eta=0.01$ and $\lambda=0.1$.
  Results are averaged across 1000 ImageNet images.
  On the left, we record the top-1 accuracy at $T{=}10$ across FGSM and PGD-10 attacks, with JPEG ($Q{=}75$) included as the non-differentiable analog.
  The parentheses show the clean--adversarial gap recovered.
  On the right, we consider values of $T$ from 0 to 30 at increments of 5 for all four attack configurations.
  }
  \label{fig:purification}
\end{figure}

Using ACE, we can examine the effects of adversarial perturbations on the DCT spectrum and how the DCT purification influences it.
Table~\ref{tab:freq} reports the mean AC DCT entropy per frequency sub-band under clean, adversarial, and purified conditions.
Adversarial perturbations inflate entropy by $+0.009$--$+0.022$ bits in \emph{every} band, not only in high frequencies as the $\ell_\infty$ attack construction might suggest.
The proxy's global AC-entropy minimization consequently reduces all bands uniformly ($-0.26$ bits total), removing adversarial content as a side effect of broadband smoothing.

\begin{table}[ht]
  \centering
  \small
  \caption{
  \textbf{Adversarial perturbations span AC coefficients uniformly, and ACE's gradient uniformly suppresses AC coefficients.}
  Mean AC DCT entropy (bits) by frequency sub-band.
  Values are averaged across 200 ImageNet images.
  }
  \label{tab:freq}
  \begin{tabular}{lcccc}
    \toprule
    Condition & Low (0--20) & Mid (21--41) & High (42--62) & Total \\
    \midrule
    Clean              & 3.418 & 2.809 & 2.676 & 2.968 \\
    FGSM ($\varepsilon=4/255$)   & 3.439 & 2.833 & 2.697 & 2.990 \\
    PGD-10 ($\varepsilon=4/255$) & 3.426 & 2.819 & 2.687 & 2.977 \\
    Purified (FGSM)    & 2.869 & 2.625 & 2.617 & 2.704 \\
    Purified (PGD)     & 2.868 & 2.624 & 2.616 & 2.703 \\
    \bottomrule
  \end{tabular}
\end{table}

Because adversarial energy is distributed broadly, a proxy restricted to high frequencies cannot remove it.
We confirm this by considering three variants on a 500-image PGD-10 evaluation.
The \emph{full} proxy over all 63 AC coefficients ($i{+}j\ge1$), a \emph{mid+high} proxy over 49 coefficients ($i{+}j\ge5$), and a \emph{high-only} proxy over 21 coefficients ($i{+}j\ge9$).

\begin{table}[ht]
  \centering
  \caption{
  \textbf{To purify adversarial perturbations, it is necessary to suppress signals across all frequency bands, not just high frequencies.}
  Here we measure ACE's ability to purify adversarial perturbations when restricted to high ($i+j\geq9$), mid-to-high ($i+j\geq5$), or full ($i+j\geq1$) AC coefficients.
  Adversarial perturbations are constructed using PGD-10 at $\varepsilon=4/255$, and $\lambda=0.05$ in \eqref{eq:purify} is used.
  }
  \label{tab:freq_ablation}
  \begin{tabular}{lcccc}
    \toprule
    ACE variant & AC coefs & $T=5$ & $T=10$ & $T=20$ \\
    \midrule
    Full & 63 & 41.4\% & 42.1\% & 24.3\% \\
    Mid-to-high & 49 & 17.4\% &  21.7\% & 18.1\% \\
    High & 21 & 0.0\% & 0.0\% &  0.0\% \\
    \bottomrule
  \end{tabular}
\end{table}

Table~\ref{tab:freq_ablation} shows that a proxy restricted to high coefficients achieves $0\%$ recovery, and a recovery restricted to mid-to-high coefficients only partially recovers accuracy. 

% ---------------------------------------------------------------------------
\subsection{Training-Time Augmentation}
\label{sec:aug}
% ---------------------------------------------------------------------------

The same ACE gradient used for inference-time purification can be applied during training.
Given a clean image $\bx$, run $T$ steps of~\eqref{eq:purify} as a form of data augmentation.
This exposes the model to images with reduced AC entropy---a structured transformation distinct from standard augmentations such as random crops, color jitter, or JPEG augmentation.

We train a ResNet-18 on CIFAR-10~\citep{krizhevsky2009learning} for 100 epochs with cosine learning rate decay, batch size 128, SGD with momentum 0.9, and weight decay $5\times10^{-4}$.
We test four augmentation pipelines: \emph{baseline} uses random crop and horizontal flip only; \emph{ACEAug} additionally considers ACEd images with $T=5$, $\lambda=0.05$, and $p=0.5$; \emph{AugMix} uses JSD consistency loss with weight 12, width 3, and severity 3 from~\citet{hendrycks2020augmix}; and \emph{AugMix+ACEAug} which additionally considers ACEAug.
All conditions use the same architecture and optimizer; results are averaged over 3 random seeds.

We report top-1 accuracy on CIFAR-10-C~\citep{hendrycks2019benchmarking}, which contains 15 corruption types, each with five severity levels.
We compute the mean accuracy over all 15 corruptions and 5 severity levels, and the mean Corruption Error (mCE) relative to AlexNet~\citep{krizhevsky2012imagenet}.

\begin{table}[h]
  \centering
  \small
  \caption{
    \textbf{ACEd images improve CIFAR-10-C corruption robustness beyond the benefits already provided by AugMix.}
    Per-category values are means over 3 seeds, with standard deviation also reported.}
  \label{tab:cifar_c}
  \begin{tabular}{lcccc}
    \toprule
    Category & Baseline & ACEAug & AugMix & AugMix+ACEAug \\
    \midrule
    Noise (gaussian, shot, impulse)             & 51.9 & 60.8 & 83.2 & \textbf{84.0} \\
    Blur (defocus, glass, motion, zoom)          & 73.5 & 79.7 & 89.9 & \textbf{90.0} \\
    Weather (snow, frost, fog, brightness)      & 85.7 & 88.5 & 91.4 & \textbf{91.4} \\
    Digital (contrast, elastic, pixelate, JPEG) & 79.2 & 84.7 & 89.5 & \textbf{89.8} \\
    \midrule
    \textbf{Mean (15 corruptions)}              & $74.0 \pm 1.3$ & $79.6 \pm 0.2$ & $88.8 \pm 0.2$ & $\mathbf{89.1 \pm 0.1}$ \\
    mCE (vs.\ AlexNet, $\downarrow$)            & $33.8 \pm 1.6$ & $26.2 \pm 0.2$ & $14.9 \pm 0.2$ & $\mathbf{14.4 \pm 0.2}$ \\
    \midrule
    Clean (CIFAR-10 test)                       & $95.0 \pm 0.2$ & $94.8 \pm 0.1$ & $95.4 \pm 0.0$ & $95.0 \pm 0.1$ \\
    \bottomrule
  \end{tabular}
\end{table}

All four corruption categories improve when ACEAug is added to the standard baseline (Table~\ref{tab:cifar_c}), with gains of $+6.2$~pp (Blur), $+5.5$~pp (Digital), $+8.9$~pp (Noise), and $+2.8$~pp (Weather).
The largest per-corruption gains are in \emph{pixelate} and \emph{zoom blur}, both of which reduce high-frequency content in a manner structurally similar to the proxy's AC-entropy minimization.
\emph{JPEG compression} also improves substantially---the direct analog of the proxy's transform-domain smoothing.

AugMix already reduces mCE to $14.9 \pm 0.2$, and adding ACEAug further reduces it by $0.5$ points, with gains visible across all four corruption categories and no meaningful clean-accuracy cost.
This additive improvement demonstrates that ACEAug captures frequency-domain structure that is orthogonal to the geometric and photometric variations covered by AugMix's augmentation chain.
ACEAug acts on AC coefficient statistics, while AugMix's operations primarily perturb spatial geometry and color, leaving complementary room for the proxy to reduce sensitivity to coding artifacts.
ACEAug mCE variance is very low ($\pm 0.2$) in both the baseline and AugMix settings, confirming that the benefit is reliable across random seeds.

% ---------------------------------------------------------------------------
\subsection{Compression-Surviving Adversarial Perturbations}
\label{sec:survival}
% ---------------------------------------------------------------------------

Section~\ref{sec:adv_pur} used ACE to \emph{remove} adversarial content by exploiting the fact that PGD perturbations inflate AC entropy.
The same signal can be reversed to craft perturbations that retain their structure even when passed through a lossy codec.
This matters for physical-world attacks (printed patches, camera-captured stickers), which must survive JPEG or HEVC transcoding at some point in the capture pipeline~\citep{shin2017jpeg}.

We add the proxy as a differentiable constraint on the perturbation during attack construction.
In particular, using PGD-10 on a ResNet-50 ImageNet classifier to minimize classifier confidence plus $\mu\,\hat{r}(\delta)$, where $\delta$ is the perturbation itself.
We measure the \emph{survival rate}, the fraction of pre-compression attack strength retained after JPEG or HEVC compression for $\mu\in\{0,1.0\}$ and $\varepsilon\in\{4/255,8/255\}$.
The clean accuracy of the classifier is $73.1\%$.

\begin{table}[h]
  \centering
  \small
  \caption{
  \textbf{ACE's proxy gradient successfully constructs adversarial perturbations to survive lossy codec compression.}
  Results are averaged across 200 images.}
  \label{tab:survival}
  \begin{tabular}{lccccc}
    \toprule
    $\varepsilon$ & $\mu$ & $\delta$-BPP & Surv.\ JPEG-50 & Surv.\ JPEG-75 & Surv.\ HEVC-32 \\
    \midrule
    \multirow{2}{*}{4/255}
      & 0 & 4.884 & 48.0\% & 77.7\% & 29.7\% \\
      & 1.0              & 4.594 & 62.4\% & 90.1\% & 34.4\% \\
      \midrule
    \multirow{2}{*}{8/255}
      & 0 & 4.830 & 71.7\% & 98.1\% & 36.8\% \\
      & 1.0              & 4.635 & 85.5\% & 97.6\% & 42.4\% \\
    \bottomrule
  \end{tabular}
\end{table}

ACE lowers the perturbation's own AC entropy ($\delta$-BPP drops $0.2$--$0.4$ bits across all $\mu$), which translates into $+14$pp survival at JPEG-50 and $+5$--$6$pp at HEVC-32 across both $\varepsilon$ (Table~\ref{tab:survival}).
Mechanistically, this is the exact mirror of Section~\ref{sec:adv_pur}: there, the proxy's gradient removes perturbations that inflate AC entropy; here, the same gradient signal is used at construction time to keep a perturbation's energy in the codec's low-entropy null space.

% ---------------------------------------------------------------------------
\subsection{Compressible Features for Split Inference}
\label{sec:split}
% ---------------------------------------------------------------------------

Split inference partitions a network's computation by transmitting the intermediate feature map over a bandwidth-constrained channel, facilitating computation across edge devices and servers~\citep{kang2017neurosurgeon}.
Existing feature-compression approaches either learn a dedicated bottleneck layer~\citep{shi2019bottlenet} or use a fixed quantization scheme; both add parameters or ignore compressibility during training.
Using ACE, we directly incorporate compressibility into the training loss by treating the intermediate feature maps as pseudo-images, $\mathcal{L} = \mathcal{L}_{\text{CE}} + \lambda\,\hat{r}(\text{features})$.

We finetune a pretrained ResNet-50 split at the layer2 output on ImageNet, SGD with frozen batch normalization, cosine learning rate, 20 epochs, comparing a baseline ($\lambda=0$) and compression-aware ($\lambda=0.1$) model.
Both reach $75.9\%$ uncompressed top-1 accuracy.
At inference time, we compress the transmitted features using JPEG at several quality levels and report top-1 accuracy as a function of bits per spatial position (bpsp).

\begin{table}[h]
  \centering
  \small
  \caption{
  \textbf{ACE's gradients can be utilized during training to improve the compression rate of intermediate features, without sacrificing performance.}
  Here, we provide a rate-distortion comparison at matched feature bitrate between the baseline and the compression-aware model.
  }
  \label{tab:split}
  \begin{tabular}{lcc}
    \toprule
    bpsp & Baseline Acc.\ & Compression-aware Acc.\ \\
    \midrule
    $\sim$1274 & $\sim$43\% & \textbf{47.7\%} \\
    $\sim$1360 & $\sim$50\% & \textbf{55.1\%} \\
    $\sim$1486 & $\sim$57\% & \textbf{61.2\%} \\
    \bottomrule
  \end{tabular}
\end{table}

The proxy shifts the feature RD curve toward the origin: at matched bpsp, the proxy-trained features give $+4$--$5$pp accuracy (Table~\ref{tab:split}), thereby reaching the same accuracy at $7$--$14\%$ lower bitrate.
This is a direct file-size effect: training with the proxy penalty reduces feature AC entropy by $0.435$ bits/position, which lowers the resulting JPEG file size by $7$--$14\%$ across all quality settings tested.
The result confirms that ACE's DCT-entropy signal is not specific to RGB pixel statistics---it transfers, with zero added parameters, to learned intermediate representations that are not natural images.

% ---------------------------------------------------------------------------
\subsection{Rate-Constrained Guidance in Diffusion Sampling}
\label{sec:diffguide}
% ---------------------------------------------------------------------------

Diffusion samplers admit differentiable guidance terms added at each denoising step, as in classifier guidance~\citep{dhariwal2021diffusion}.
Adding $\lambda\,\nabla_{\bx_t}\hat{r}(\hat{\bx}_0)$ to the DDIM~\citep{song2021ddim} update steers generation toward images that are cheaper to code---useful for streaming or broadcast pipelines where generated content must not cause bitrate spikes.
Unlike classifier or CLIP guidance~\citep{mishra2024image}, ACE requires no auxiliary network.

We guide Stable Diffusion~1.5~\citep{rombach2022latent} (using 50 DDIM steps, with guidance applied only on the last 25 steps), and measure mean bitrate, FID~\citep{heusel2017fid}, and CLIP score~\citep{hessel2021clipscore}.
We sweep $\lambda\in\{0,20,50,100,200,400,800\}$, to trace the full fidelity--rate trade-off.

\begin{figure}[ht]
  \centering
  \begin{tikzpicture}
  \begin{groupplot}[
    group style={group size=3 by 1, horizontal sep=1.4cm},
    width=0.34\linewidth,
    height=0.28\linewidth,
    xmode=log,
    log basis x={2},
    xlabel={$\lambda$},
    xmin=15, xmax=1000,
    tick label style={font=\scriptsize},
    label style={font=\footnotesize},
    legend style={font=\scriptsize},
  ]
  \nextgroupplot[
    ylabel={Bitrate},
  ]
    \addplot[color=black!50, dashed, line width=1pt, domain=15:2000] {0.4194};
    \addplot[color=proxyblue, line width=1.8pt, mark=*, mark size=1.6pt,
             mark options={fill=proxyblue}]
      table[col sep=comma, x=lam, y=hevc_bpp] {figures/data/diffusion_guidance_v3.csv};
  \nextgroupplot[
    ylabel={FID},
  ]
    \addplot[color=black!50, dashed, line width=1pt, domain=15:2000] {24.26};
    \addplot[color=proxyblue, line width=1.8pt, mark=*, mark size=1.6pt,
             mark options={fill=proxyblue}]
      table[col sep=comma, x=lam, y=fid] {figures/data/diffusion_guidance_v3.csv};
  \nextgroupplot[
    ylabel={CLIP score},
  ]
    \addplot[color=black!50, dashed, line width=1pt, domain=15:2000] {31.24};
    \addplot[color=proxyblue, line width=1.8pt, mark=*, mark size=1.6pt,
             mark options={fill=proxyblue}]
      table[col sep=comma, x=lam, y=clip_score] {figures/data/diffusion_guidance_v3.csv};
  \end{groupplot}
  \end{tikzpicture}
  \caption{
  \textbf{ACE-guided image generation meaningfully improves the compressibility of outputs whilst maintaining their overall quality}.
  Bitrate is measured using HEVC CRF~32.
  The dashed line gives the unguided $\lambda$ equal to $0$ baseline.
  Values are measured against 5000 COCO-val references~\citep{lin2014coco}.
  }
  \label{fig:diffguide}
\end{figure}

Figure~\ref{fig:diffguide} shows that bitrate and FID both change monotonically with $\lambda$, but at different rates, allowing for favorable trade-offs.
Up to $\lambda{=}100$, bitrate falls steeply---down $41\%$ relative to the baseline---while FID rises only by $+3.9\%$ and CLIP score is unchanged.
In this regime, the sampler is finding semantically equivalent, more-compressible images at essentially no perceptual cost.
Beyond $\lambda{=}100$ the trade-off steepens, with the guidance increasingly trading away image fidelity for compressibility.